\documentclass[]{article}
\usepackage{graphicx} % Required for inserting images
\usepackage{amsmath}
\usepackage{amssymb}
\usepackage{amsfonts}
\usepackage{stmaryrd}
\usepackage{titlesec}
\usepackage[margin=1in]{geometry}

\setlength{\parindent}{0pt}

\newcommand{\msp}{\text{ }}
\newcommand{\ligneh}{\msp | \msp}
\newcommand{\si}{\text{ si }}
\newcommand{\et}{\text{ et }}
\newcommand{\ou}{\text{ ou }}
\newcommand{\ch}{\text{ch}}
\newcommand{\sh}{\text{sh}}
\newcommand{\rg}{\text{rg}}
\newcommand{\Vect}{\text{Vect}}
\newcommand{\dl}[2]{\text{DL}_{#1}(#2)}
\newcommand{\mat}{\text{Mat}}
\newcommand{\ind}{\perp\!\!\!\perp} 
\newcommand{\conv}{\text{ converge }}
\newcommand{\cov}{\text{Cov}}
\newcommand{\tr}{\text{tr}}

\newcommand{\dd}{\text{d}}
\newcommand{\dx}{\text{d}x}
\newcommand{\dy}{\text{d}y}
\newcommand{\dz}{\text{d}z}
\newcommand{\dr}{\text{d}r}
\newcommand{\dt}{\text{d}t}
\newcommand{\dtheta}{\text{d}\theta}
\newcommand{\dphi}{\text{d}\phi}
\newcommand{\repr}{|\mathcal{R}}
\newcommand{\grad}{\overrightarrow{{\text{grad}}} \;}
\newcommand{\cste}{\text{cste}}

\setlength{\emergencystretch}{3em}
\hbadness=10000

\title{Demonstration de cours semaine 1}
\author{Huang Ziwen}
\date{2026-2027 PC*}

\begin{document}

\section{Théorème du rang (forme géométrique)}

Soit $E \et F$ des K-evs, $u \in \mathcal{L}(E, F)$ et $E = \ker u \oplus E'$, alors \underline{$u$ induit un isomorphisme de $E'$ sur $\text{Im } u$} \\

Soit $\tilde{u} : E' \rightarrow \text{Im } u, x \mapsto u(x)$. \\

$\tilde{u}$ est bien défini et linéaire. 

\begin{itemize}
    \item $\ker \tilde{u} = \ker u \cap E' = \{0\}$ donc $\tilde{u}$ est injective
    \item Soit $y \in \text{Im}(u)$. Alors $\exists x \in E$ tq $u(x) = y$. \\
    
    Comme $E = \ker u \oplus E'$, fixons $x_0 \in \ker u, x' \in E'$ tq $x = x_0 + x'$. 

    \[ y = u(x) = u(x_0 + x') = u(x_0) + u(x') = 0 + u(x') = \tilde{u}(x') \]

    Donc $\tilde{u}$ est surjective
\end{itemize}

Donc $u$ induit un isomorphisme de $E'$ sur $\text{Im } u$. 

\section{Définition d'un projecteur}

Soit $p \in \mathcal{L}(E)$. $p$ est un projecteur si \underline{$p^2 = p$}. \\

\underline{$E = \ker p \oplus \text{Im }p$.}

\begin{itemize}
    \item Montrons $E = \ker p + \text{Im }p$ \\
    
    On a $x = p(x) + (x - p(x))$. $p(x) \in \text{Im }p$. \\

    $p( x - p(x) ) = p(x) - p^2(x) = p(x) - p(x) = 0$, donc $(x-p(x)) \in \ker p$. \\

    Donc $E = \ker p + \text{Im }p$
    \item Montrons que la somme est directe. \\
    
    Soit $y \in \ker p \cap \text{Im }p$. Alors $\exists x \in E$ tq $y = p(x)$.

    \[ 0 = p(y) = p^2(x) = p(x) = y \]

    Donc $y = 0$ d'où $\ker p \cap \text{Im }p \subset \{0\}$. La somme est directe
\end{itemize}

\underline{$s = 2p - \text{Id}$ est une symétrie.} 

\[ s^2 = (2p - \text{Id})\circ(2p - \text{Id}) = 4p^2 - 2p - 2p + \text{Id} = 4p - 4p + \text{Id} = \text{Id} \]

\section{Définition d'une somme directe}

Soit $(E_i)_{i \in I}$ une famille finie de sous-espaces vectoriels de $E$. Les $E_i$ sont en somme directe si \underline{$\forall x \in \sum_{i \in I} E_i$, la décomposition $x = \sum_{i \in I} x_i$ avec $(x_i)_{i \in I} \in \prod_{i \in I} E_i$ est unique}. \\

\underline{Les $E_i$ sont en somme directe $\iff \forall (x_i)_{i \in I} \in \prod_{i \in I} E_i, (\sum_{i \in I} x_i = 0) \implies (\forall i \in I, x_i = 0)$}. \\

Posons 

\begin{align*}
    \Phi: \prod_{i \in I} E_i \rightarrow & \sum_{i \in I} E_i \\
    (x_i)_{i \in I} \mapsto & \sum_{i \in I} x_i
\end{align*}

$\Phi$ est linéaire et surjective.

\[ \Phi \text{ est un isomorphisme} \iff \Phi \text{ est injective} \iff \forall x \in \sum_{i \in I} E_i, \exists ! (x_i)_{i \in I} \in \prod_{i \in I} E_i \ligneh x = \sum_{i \in I} x_i\]

Donc les $E_i$ sont en somme directe $\iff \Phi $ est un isomorphisme. Or, $\Phi $ est un isomorphisme $\iff \Phi$ est injective $\iff \ker \Phi = \{0\}$. D'où

\[ \text{Les } E_i \text{ sont en somme directe} \iff \forall (x_i)_{i \in I} \in \prod_{i \in I} E_i, (\sum_{i \in I} x_i = 0) \implies (\forall i \in I, x_i = 0) \]

\underline{$\dim(\sum_{i \in I}E_i) \leq \sum_{i \in I} \dim E_i$} \\

Comme $\Phi$ est surjective, $\dim(\sum_{i \in I}E_i) = \rg \Phi$. \\

Par théorème du rang, $\rg \Phi \leq \dim (\prod_{i \in I} E_i) = \sum_{i \in I} \dim E_i$. D'où $\dim(\sum_{i \in I}E_i) \leq \sum_{i \in I} \dim E_i$ \\

\underline{$\dim(\sum_{i \in I}E_i) = \sum_{i \in I} \dim E_i \iff $ les $E_i$ sont en somme directe} \\

\[ \dim(\sum_{i \in I}E_i) = \sum_{i \in I} \dim E_i \iff \Phi \text{ est un isomorphisme} \iff \text{les } E_i \text{ sont en somme directe} \]

\section{$\ker u \et \text{Im } u$ stables par $v$}

Soit $(u, v) \in \mathcal{L}(E)^2$ tq $u \et v$ commutent. \\

\underline{$\ker u$ est stable par $v$.} \\

Soit $x \in \ker u$. $u(x) = 0$

\[ 0 = v(u(x)) = u(v(x)) \]

Donc $v(x) \in \ker u$ \\

\underline{$\text{Im } u$ est stable par $v$.} \\

Soit $y \in \text{Im } u$. Fixons $x \in E$ tq $y = u(x)$.

\[ v(y) = v(u(x)) = u(v(x)) \in \text{Im } u \]

Donc $v(y) \in \text{Im } u$ \\

En particulier, $u$ commute avec $P(u)$ donc \underline{$\ker P(u) \et \text{Im } P(u)$ sont stables par $u$}

\section{$F$ un ssev stable par $u$}

Soit $E$ un K-ev de dimension $n$, $u \in \mathcal{L}(E)$, $F$ un ssev de $E$ de dimension $p$ stable par $u$. Soit $\mathcal{B} = (e_i)_{i \in \llbracket 1, n \rrbracket}$ une base adaptée à $F$. 

\[ \mat_{\mathcal{B}}(u) = \begin{pmatrix}
    A & B \\
    (0) & C
\end{pmatrix} \]

En particulier, $A = \mat_{\mathcal{B}'}(u_F)$ avec $\mathcal{B}' = (e_i)_{i \in \llbracket 1, p \rrbracket}$ \\

\underline{$ \begin{vmatrix}
    A & B \\
    (0) & C
\end{vmatrix} = \det A \times \det C$}

\[ \begin{pmatrix}
    A & B \\
    (0) & C
\end{pmatrix} = \begin{pmatrix}
    I_p & (0) \\
    (0) & C
\end{pmatrix} \times \begin{pmatrix}
    A & B \\
    (0) & I_q
\end{pmatrix} \]

$\begin{vmatrix}
    I_p & (0) \\
    (0) & C
\end{vmatrix} = \det C$ en développant par rapport aux $p$ premières colonnes. \\

$\begin{vmatrix}
    A & B \\
    (0) & I_q
\end{vmatrix} = \det A$ en développant par rapport aux $q$ dernières lignes. \\

D'où $ \begin{vmatrix}
    A & B \\
    (0) & C
\end{vmatrix} = \det A \times \det C$

\section{Trace : définitions et propriétés}

\begin{align*} 
    \tr : \mathcal{M}_n(\mathbb{K}) \rightarrow & \mathbb{K} \\
    A \mapsto & \sum_{i=1}^n [A]_{i, i} 
\end{align*}

La trace est une \underline{forme linéaire}. \\

Soit $A = (a_{i, j})_{(i, j) \in \llbracket 1, n \rrbracket \times \llbracket 1, p \rrbracket} \in \mathcal{M}_{n, p}(\mathbb{K}), B = (b_{i, j})_{(i, j) \in \llbracket 1, p \rrbracket \times \llbracket 1, n \rrbracket} \in \mathcal{M}_{p, n}(\mathbb{K})$. \\

\underline{$\tr (AB) = \tr (BA)$} \\

\[ \tr (AB) = \sum_{i = 1}^{n} (\sum_{j = 1}^{p} a_{i, j} b_{j, i}) = \sum_{j = 1}^{p} \sum_{i = 1}^{n} b_{j, i} a_{i, j} = \tr(BA) \]

Soit $A \in \mathcal{M}_n(\mathbb{K}), P \in \text{GL}_n(\mathbb{K})$. \underline{$\tr(P^{-1}AP) = \tr A$}

\[ \tr(P^{-1}AP) = \tr((P^{-1}A)P) = \tr(PP^{-1}A) = \tr A \]

\section{Polynômes de Lagrange}

Soit $n \in \mathbb{N}^*, a = (a_i)_{i \in \llbracket 1, n \rrbracket} \in \mathbb{K}^n$. \underline{$\forall i \in \llbracket 1, n \rrbracket, L_i$ est l'unique polynôme $\in \mathbb{K}_{n-1}[X]$ tel que $L_i(a_j) = \delta_{i, j}$}.

\[ L_i = \frac{\prod_{j=1, j\neq i}^{n} (X-a_j)}{\prod_{j=1, j\neq i}^{n} (a_i-a_j)} \]

\underline{$(L_i)_{i \in \llbracket 1, n \rrbracket}$ est une base de $\mathbb{K}_{n-1}[X]$}.\\

Soit $(\lambda_1, \dots , \lambda_n) \in \mathbb{K}^n$ tq $\sum_{i=1}^n \lambda_i L_i = 0$. En particulier, $\forall k \in \nint{1}{n}$, 

\[ \sum_{i=1}^{n}\lambda_i L_i(a_k) = \sum_{i=1}^{n}\lambda_i \delta_{i, k} = \lambda_k \]

Donc $\forall k \in \nint{1}{n}, \lambda_k = 0$. \\

$(L_i)_{i \in \llbracket 1, n \rrbracket}$ est libre maximale dans $\mathbb{K}_{n-1}[X]$, donc elle est une base. \\

\underline{$\forall P \in \mathbb{K}_{n-1}[X], P = \sum_{i = 1}^{n} P(a_i)L_i$}. \\

Soit $P \in \mathbb{K}_{n-1}[X]$. Notons $(\lambda_i)_{i \in \nint{1}{n}}$ ses coordonnées dans la base $(L_i)_{i \in \llbracket 1, n \rrbracket}$. \\

$\forall k \in \nint{1}{n}$,

\[ P(a_k) = \sum_{i=1}^{n} \lambda_i L_i(a_k) = \sum_{i=1}^{n} \lambda_i \delta_{i, k} = \lambda_k \]

Donc $\forall k \in \nint{1}{n}$, $\lambda_k = P(a_k)$. Alors $P = \sum_{i = 1}^{n} P(a_i)L_i$

\section{Déterminant de Vandermonde}

Soit $n \in \mathbb{N}^*, (a_1, \dots, a_n) \in \mathbb{K}^n$. 

\[ V(a_1, \dots, a_n) = \begin{vmatrix}
    1 & a_1 & \dots & a_1^{n-1} \\
    1 & a_2 & \dots & a_2^{n-1} \\
    \vdots & \vdots &  & \vdots \\
    1 & a_n & \dots & a_n^{n-1} \\
\end{vmatrix} \]

\underline{$V(a_1, \dots, a_n) = \prod_{1 \leq i < j \leq n} (a_j - a_i)$} \\

Procédons par récurrence sur $n$. $\forall n \in \mathbb{N}^*$, $H_n$:``$\forall (a_1, \dots, a_n) \in \mathbb{K}^n, V(a_1, \dots, a_n) = \prod_{1 \leq i < j \leq n} (a_j - a_i)$''.

\begin{itemize}
    \item Initialisation : pour $n = 1$, $\begin{vmatrix}
        1
    \end{vmatrix} = 1$
    \item Hérédité : soit $n \in \mathbb{N}^*$ tq $H_n$. Montrons $H_{n+1}$ \\
    
    Soit $(a_1, \dots, a_{n+1}) \in \mathbb{K}^{n+1}$
    \begin{itemize}
        \item Si $\exists (p, q) \in \nint{1}{n+1}^2$ tq $p \neq q \et a_p = a_q$, la matrice possède 2 lignes identiques. \\
        
        Alors $V(a_1, \dots, a_{n+1}) = 0 = \prod_{1 \leq i < j \leq n} (a_j - a_i)$ \\

        \item Sinon, les $(a_1, \dots, a_{n+1})$ sont 2 à 2 distincts. \\
        
        Considérons $f:x \mapsto V(a_1, \dots, a_n, x) = \begin{vmatrix}
            1 & a_1 & \dots & a_1^n \\
            1 & a_2 & \dots & a_2^n \\
            \vdots & \vdots &  & \vdots \\
            1 & a_n & \dots & a_n^n \\
            1 & x & \dots & x^n \\
        \end{vmatrix}$ \\

        En développant suivant la dernière ligne, $f$ est une fonction polynomiale de degré $\leq n$. Son coefficient dominant $\alpha = \begin{vmatrix}
            1 & a_1 & \dots & a_1^{n-1} \\
            1 & a_2 & \dots & a_2^{n-1} \\
            \vdots & \vdots &  & \vdots \\
            1 & a_n & \dots & a_n^{n-1} \\
        \end{vmatrix} = V(a_1, \dots, a_n)$. \\

        Or, $\forall k \in \nint{1}{n}, f(a_k) = 0$ car la matrice aurait 2 lignes identiques. On a donc toutes les $n$ racines de $f$. \\

        Donc $\forall x \in \mathbb{K}, f(x) = \alpha \prod_{i=1}^{n}(x-a_i)$. En particulier, 

        \begin{align*}
            V(a_1, \dots, a_{n+1}) &= f(a_{n+1}) \\
            &= V(a_1, \dots, a_n) \prod_{k=1}^{n}(a_{n+1}-a_k) \\
            &= \prod_{1 \leq i < j \leq n} \prod_{k=1}^{n} (a_j - a_i) (a_{n+1}-a_k) \\ 
            &= \prod_{1 \leq i < j \leq n+1} (a_j - a_i)
        \end{align*} 
    \end{itemize}
\end{itemize}

\section{Polynômes d'endomorphismes et de matrices}

Soit $E$ un K-ev, $u \in \mathcal{L}(E), P = \sum_{k} a_k X^k \in \mathbb{K}[X]$. \\

\underline{$P(u) = \sum_{k} a_k u^k$} \\

$P \mapsto P(u)$ est \underline{linéaire} et compatible avec la composition \underline{$\forall (P, Q) \in \mathbb{K}[X]^2, (PQ)(u) = P(u) \circ Q(u)$} (morphisme d'algèbre). \\

Soit $Q = \sum_{i} b_i X^i \in \mathbb{K}[X], k \in \mathbb{N}$.

\begin{align*}
    X^kQ(u) &= (\sum_{i} b_i X^{i+k}(u)) \\
    &= \sum_{i} b_i u^{i+k} \\
    &= u^k \circ (\sum_{i} b_i u^i)
\end{align*}

Soit $P = \sum_{k} a_k X^k \in \mathbb{K}[X]$. Alors $PQ = \sum_{k} a_k X^kQ$

\[ PQ(u) = \sum_{k}a_k u^k \circ Q(u) = P(u) \circ Q(u) \]

\end{document}